\makeabstract
{
Effects of Swing Team Routing Policy on Emergency Department Performance
}
{
Emile Ndayishimiye(1), Mayisa Tasnum(1), Professor. Kristy Gardner(1)
}
{
\textsuperscript{1}Department of Computer Science, Amherst College, Amherst, MA, USA
}
{
Emergency Department (ED) crowding remains a critical issue, with prolonged length of stay (LOS) negatively impacting patient outcomes. This study applies queueing theory to evaluate policies for routing patients to a flexible 'Swing team' of providers, with the objective of reducing LOS for moderate-acuity patients. 
}
{
In this study, we use a simulation modeled on patient flow at Mt. Sinai Hospital ED. We tested four distinct policies to assign patients from triage to the treatment zone. These included Join-the-Shortest-Queue (JSQ), which routes patients to any zone with the shortest line, and Swing Extends Red (SER), where the Swing team assists the high-moderate acuity 'Red zone’. We also tested two load-balancing policies: Idealized Load Balancing (ILB), which uses staffing levels and queue lengths to route patients to the zone with the shortest expected wait time, and Constrained Load Balancing (CLB), which first directs all high-moderate acuity patients to the Red zone before applying the ILB logic for other patients. 
}
{
Results indicate that ILB achieves the lowest LOS overall (a 58.9\% reduction from baseline with no Swing team) but violates operational constraints by routing high-moderate acuity patients to zones other than Red. In contrast, CLB balances workload while adhering to the zone constraints, reducing LOS by a substantial 31.5\%. SER performs best when swing staffing is low, while ILB and CLB excel under higher swing capacity. 
}
{
In conclusion, workload balancing policies show significant promise for improving ED efficiency. The CLB policy, in particular, emerges as the optimal and implementable strategy. It effectively balances the high performance of an idealized model with the practical, real-world constraints of patient care, significantly reducing LOS. These findings demonstrate that routing policies are a powerful tool to effectively utilize the Swing team's potential to contribute to ED performance.
}

\newpage
